In Section \ref{sec:measurementofdnps}, the various methods for experimentally understanding delayed neutron precursors were discussed.
The next stage of understanding these precursors is to understand how they are used in kinetics models, as well as the various ways they are represented.
The microscopic and macroscopic experimental approaches provide the necessary data to create more easily handled forms.
By summing the contributions from each \gls{DNP} in the microscopic approach, the macroscopic and microscopic approaches can be said to provide data in the form of a delayed neutron count rate over time.
Section \ref{sec:calc-group-param} discusses the least-squares approach to convert this data into several \gls{DNP} groups.

The \glspl{DNP} are commonly represented in six groups.
These groups were developed originally for computational efficiency and due to lacking data for individual precursors.
Brady and England found that increasing the number of groups for point kinetics calculations from six to nine caused no significant differences \cite{brady1989delayed}.
Their study shows that increasing the number of \gls{DNP} groups alone does not provide an increase in accuracy for point kinetics simulations.

Representing the \glspl{DNP} in group form also simplifies the calculation of various kinetics parameters.
Various kinetics parameters are introduced in Section \ref{sec:bg-trans-analysis-methods}, though without detail on the purpose of each variable.
This section provides more information on what these kinetics and delayed neutron group parameters physically represent, while Sections \ref{sec:calc-group-param} and \ref{sec:calc-kinetic-param} detail how they are calculated.
The \gls{DNP} group parameters directly affect many of the kinetics parameters, and both sets of data are used in kinetics and dynamics problems.
Table \ref{tab:dnpgandk} goes through each of the terms of importance for this section.

\begin{table}[H]
    \centering
    \begin{tabular}{l l}
        \hline
        Parameter & Description \\
        \hline
        $\lambda_k$/$\lambda_i$ & The decay constant for \gls{DNP} group $k$ / nuclide $i$\\
        $\bar{\nu}_{d, k}$/$\bar{\nu}_{d, i}$ & The delayed neutron yield for \gls{DNP} group $k$ / nuclide $i$\\
        $\chi_k$/$\chi_i$ & The spectra of \gls{DNP} group $k$ / nuclide $i$\\
        $P_{n, i}$ & The delayed neutron emission probability for nuclide $i$\\
        $< T >$ & The mean half-life\\
        $Y$ & The total delayed neutron yield\\
        \hline
        $\beta_k$/$\beta_i$ & The delayed neutron fraction of \gls{DNP} group $k$ / nuclide $i$\\
        $\beta$ & The total delayed neutron fraction\\
        $\beta_{eff}$ & The effective delayed neutron fraction\\
        $\Lambda$ & The mean neutron generation time\\
        $\Lambda_{eff}$ & The effective mean neutron generation time\\
        \hline
    \end{tabular}
    \caption{\Gls{DNP} group parameters and kinetics parameters along with their definitions.}
    \label{tab:dnpgandk}
\end{table}

The \gls{DNP} groups come with associated parameters, such as the group decay constants, $\lambda_k$, and the group yields, $\bar{\nu}_{d, k}$.% (or $\beta_i$).
These group yields can be converted to group delayed neutron fractions, $\beta_k$, by dividing the group yields by the average number of neutrons emitted per fission event.
The group yields represent the number of delayed neutrons produced per fission in group $i$, and when summed yield the total delayed neutron yield per fission.
The $\beta_k$ and $\lambda_k$ terms are the defining parameters for the \gls{DNP} groups, along with the group spectra, $\chi_k$, for \gls{DNP} group $k$.
With this data, the other group parameters can be derived, such as the total delayed neutron yield, $Y$, and delayed neutron fraction $\beta$.

Although the groups are useful, there are times when individual \glspl{DNP} are needed.
For example, in 1986, Perry et al. demonstrates that kinetics calculations yield different results when using individual nuclides compared to the six-group parameters \cite{perry1986application}.
In 1989, Brady and England also use the individual nuclides to compare with group data \cite{brady1989delayed}, and  updated the six-group parameters based on the explicit \glspl{DNP}.
Updating the six-group parameters enabled greater consistency between experimental results and simulations.

%Representing the \glspl{DNP} in group form also simplifies the calculation of various kinetics parameters.
%Various kinetics parameters are introduced in Section \ref{sec:bg-trans-analysis-methods}, though without detail on the purpose of each variable.
%This section provides more information on what these kinetics and delayed neutron group parameters physically represent, while Sections \ref{ref:calc-group-param} and \ref{ref:calc-kinetic-param} detail how they are calculated.

Moving from the group parameters to the kinetics parameters, there is a convenient analogy which can be used to better understand some of the kinetics parameters.
The effective multiplication factor, $k_{eff}$, is an important term in reactor physics.
This term represents the scaling of the number of neutrons over each generation in a finite reactor system.
The infinite multiplication factor, $k_{\infty}$, represents the multiplication without considering leakage terms.
This is analogous to the delayed neutron fraction and mean generation time terms, where the effective form of the term requires information on the specific reactor geometry.

The adjoint flux is important for calculating the effective form of many kinetics parameters.
The adjoint flux is not a physical quantity, but rather serves as an importance, or weighting, function, describing the importance of the neutron flux in the reactor.
For example, areas further from the center of the core are less important than those closer to the center, which is reflected by a decrease in the effective delayed neutron fraction in \glspl{MSR} as \glspl{DNP} move away from the center of the core due to fuel salt advection.
The adjoint flux can be calculated by taking the adjoint of the Boltzmann transport equation, or by using the \gls{IFP} method, which provides an importance term proportional to the adjoint flux \cite{terranova2018adjoint, nauchi2010development}.

The delayed neutron fraction, $\beta$, represents the ratio of delayed neutron production over all neutron production \cite{meulekamp2006calculating}.
Related to the delayed neutron fraction is the effective delayed neutron fraction, $\beta_{eff}$.
This term is more sophisticated, as rather than being a physical property of a fissile nuclide, it depends on the composition and geometry of the reactor.
For \glspl{MSR}, it also accounts for the movement of the \gls{DNP} in different parts of the reactor.

A straightforward way to interpret the effective delayed neutron fraction is that it is the delayed neutron fraction corrected for the actual composition and geometry of the reactor.
More specifically, the effective delayed neutron fraction represents the ratio between the delayed-neutron induced fission to the total number of fission events weighted by the spectrum and adjoint neutron flux \cite{romero2022calculation}.
Alternatively, it can be thought of as the ratio of the spectrum weighted and adjoint weighted delayed neutron production over all spectrum weighted and adjoint weighted neutron production \cite{meulekamp2006calculating}.
For \glspl{MSR}, the movement of the fuel salt also affects the effective delayed neutron fraction, as the integrated adjoint neutron flux is reduced due to movement of the \glspl{DNP} out of the core and into less neutronically-important regions.

Another term, the mean neutron generation time, $\Lambda$, is the mean time between generation of a neutron and fission induced by that neutron \cite{romero2022calculation}.
This term also has an effective form, $\Lambda_{eff}$, which is the adjoint weighted form.
The effective mean neutron generation time can vary significantly, on the order of hundreds of percent, from the mean neutron generation time \cite{verboomen2006monte}.
In contrast, the effective delayed neutron fraction varies by on the order of 20\% from the delayed neutron fraction \cite{verboomen2006monte}.

\subsubsection{Calculation of DNP Group Parameters}
\label{sec:calc-group-param}

The \gls{DNP} group parameters are important for simulating transients.
The \gls{DNP} group parameters are formulated based on the fitting of the delayed neutron count rate, as discussed in Section \ref{sec:measurementofdnps}.
Equation \eqref{eq:mdnp-dn} shows the exact formulation for the delayed neutron count rate as a function of time measured by a detector for $I$ individual \glspl{DNP} for any number of fissile nuclides, recreated here as:
\begin{equation}
\dot{n}_d(t) = \epsilon \sum_{i=1}^{I} P_{n, i} \lambda_i N_{i}(t).
\nonumber
\end{equation}
Equation \eqref{eq:delnu-fit} instead gives the $K$ group \gls{DNP} fit measured by the same detector for a pulse irradiation as:
\begin{equation}
\dot{n}_d(t) \approx \epsilon F_s \sum_{k=1}^{K}  \bar{\nu}_{d, k} \lambda_k e^{-\lambda_k t}.
\label{eq:delnu-fit}
\end{equation}
For an infinite irradiation (also referred to as a saturation irradiation), Equation \eqref{eq:delnu-fit-sat} is used rather than Equation \eqref{eq:delnu-fit}, given as \cite{brady1988evaluation}:
\begin{equation}
\dot{n}_d(t) \approx \epsilon \dot{F}_s \sum_{k=1}^{K}  \bar{\nu}_{d, k} e^{-\lambda_k t}.
\label{eq:delnu-fit-sat}
\end{equation}
The terms in Equations \eqref{eq:delnu-fit} and \eqref{eq:delnu-fit-sat} are defined as follows:
\begin{align}
K &= \mbox{ the number of \gls{DNP} groups [-]} \nonumber \\ 
\lambda_k &= \mbox{ the decay constant of \gls{DNP} group \textit{k} [$s^{-1}$] } \nonumber \\ 
F_s &= \mbox{ the total number of fission events in the pulse irradiation [-]} \nonumber \\ 
\dot{F}_s &= \mbox{ the fission rate for a saturation irradiation [$s^{-1}$]} \nonumber \\ 
\bar{\nu}_{d, k} &= \mbox{ the group yield for group \textit{k} [-] .} \nonumber
\end{align}
More information on the derivation of these equations can be found in Appendix \ref{ref:appendixA}.
These equations are specifically for a single fissile nuclide, and the derivation can be found in the work of Wilson and England \cite{wilson2002delayed}.
%To make this equation generic, the $\bar{\nu}_{d, k} \lambda_k$ term and $\bar{\nu}_{d, k}$ term can instead be written as $A_k$.
%Additionally, the fission term can be written as $\tilde{F}_s$, to represent either the fission rate or the net number of fissions.
%These simplified equations do not account for decay chain daughter products \cite{conant1971measurement}.







%\begin{description}
%    %\setlength\itemsep{-0.25em}
%    \item $I$ is the total number of \glspl{DNP}
%    \item $n$ is the number of \gls{DNP} groups
%    \item $\epsilon$ is the detector efficiency
%    \item $P_{n, i}$ is the emission probability of \gls{DNP} $i$
%    \item $\lambda_i$ is the decay constant of the $i$th \gls{DNP} or group
%    \item $N_i(t)$ is the time dependent number of atoms of nuclide $i$
%    \item $F_s$ is the total number of fission events in the pulse irradiation
%    \item $\dot{F}_s$ is the fission rate for a saturation irradiation
%    \item $a_i$ is the $i$th group yield, or group abundance
%\end{description}

The pulse irradiation approach gives better statistics on the short lives \glspl{DNP}, while the longer lived \glspl{DNP} are better represented by saturation irradiation.
Brady discusses this, and uses both pulse and saturation irradiation to form groups \cite{brady1988evaluation}.
Specifically, Brady uses the four long-period groups from the saturation irradiation and the four short-period groups from the pulse irradiation renormalized at the second \gls{DNP} group.
%"...determination of four long-period groups from the infinite irradiation data and
%four short-period groups from the prompt [pulse] data, later renormalizing the two sets
%of yields at the second delayed neutron group."

%Another benefit of including saturation irradiation data is that it can be used to get the average number of neutrons, prompt and delayed, per fission \cite{conant1971measurement}.
%This is shown in Equation \eqref{eq:delnu-init}.

%\begin{equation}
%\dot{n}_d(0) = \epsilon \dot{F}_s \bar{\nu}
%\label{eq:delnu-init}
%\end{equation}


Once the delayed neutron count rate is acquired, an over-determined non-linear least squares approach is used to calculate the group yields and decays.
For example, Brady used STEPIT to handle the non-linear least-squares solve with both pulse and infinite irradiation experimental data \cite{brady1988evaluation}.
Waldo et al. also solve the least-squares problem, though the specific approach used is not specified \cite{waldo1981delayed}.
The general form of the solve is given as:
\begin{equation}
A \vec{x} \approx \vec{b}.
\label{eq:lstsq}
\end{equation}
\begin{comment}
Equation \eqref{eq:minthis} gives the $\chi^2$ term the non-linear least squares minimizes as:
\begin{equation}
    \chi^2 = \sum_{m=0}^{M} \left(\frac{\vec{b}^{(m)} - \left(A\vec{x}\right)^{(m)}}{\vec{\sigma}^{(m)}} \right) ^2,
    \label{eq:minthis}
\end{equation}
where,
\begin{align}
    \chi^2 =& \text{ the chi-squared measure of the fit } \nonumber \\
    M =& \text{ the total number of time steps over which the sample decays } \nonumber \\
    \sigma =& \text{ the uncertainty at time step $m$, or unity if the uncertainty is not measured.} \nonumber
\end{align}
\end{comment}

Based on the irradiation type used, the coefficient matrix will change, as shown in Equations \eqref{eq:delnu-fit} and \eqref{eq:delnu-fit-sat}. 
Equation \eqref{eq:lambda} gives the coefficient matrix for a pulse irradiation as:

\begin{equation}
A = \epsilon F_s
\begin{bmatrix}
\lambda_1 e^{-\lambda_1 t_0} & \lambda_2 e^{-\lambda_2 t_0} & \hdots &  \lambda_n e^{-\lambda_n t_0}\\
\lambda_1 e^{-\lambda_1 t_1} & \lambda_2 e^{-\lambda_2 t_1} & \hdots &  \lambda_n e^{-\lambda_n t_1}\\
\vdots                       & \vdots                       & \ddots &  \vdots                      \\
\lambda_1 e^{-\lambda_1 t_m} & \lambda_2 e^{-\lambda_2 t_m} &  \hdots  & \lambda_n e^{-\lambda_n t_m}\\
\end{bmatrix},
\label{eq:lambda}
\end{equation}
though for a saturation irradiation the only change would be removing the $\lambda_i$ term.

Within this least squares equation, the solution vector contains only the group yields, given in Equation \eqref{eq:lstsq-solvec} as:
\begin{equation}
\vec{x} = 
\begin{bmatrix}
\bar{\nu}_{d, 1}\\
\bar{\nu}_{d, 2}\\
\vdots\\
\bar{\nu}_{d, K}\\
\end{bmatrix}.
\label{eq:lstsq-solvec}
\end{equation}

The target vector contains the linearized delayed neutron count rate data normalized to the number of fissions, or fission rate, and the detector efficiency, as shown in Equation \eqref{eq:lstsq-target} as:
\begin{equation}
\vec{b} = 
\begin{bmatrix}
\ln \left( n_d(t_0)  \right)\\
\ln \left( n_d(t_1) \right)\\
\vdots\\
\ln \left( n_d(t_m) \right)\\
\end{bmatrix}.
\label{eq:lstsq-target}
\end{equation}
The natural logarithm is needed here, as otherwise the short time steps will dominate the generated group parameters.
Another way to resolve this is to use a relative residual, as given by Equation \eqref{eq:residual}.
\begin{equation}
    r = min_x \left| \left| \frac{\vec{b} - A\vec{x}}{\vec{b} + \epsilon} \right| \right|,
    \label{eq:residual}
\end{equation}
where,
\begin{align}
    \epsilon &= \text{small value to avoid division by zero} \nonumber \\
    &= 1\times10^{-12}. \nonumber
\end{align}

By solving this non-linear system of equations, the group half-lives, $T_k$, and yields, $\bar{\nu}_{d, k}$, can be generated.
Once this group data is generated, and if there are multiple fissioning nuclides, $F$, then Equation \eqref{eq:lami-multinuc} can be used to determine the $K$ group half lives using \cite{das1993importance} as:

\begin{equation}
\lambda_k = \frac{\sum_{F} \lambda_k^{(F)} \beta_{k, eff}^{(F)}}{\sum_F \beta_{k, eff}^{(F)}},
\label{eq:lami-multinuc}
\end{equation}
where,
\begin{align}
    \lambda_k^{(F)} =& \text{ the decay constant of \gls{DNP} group $k$ from fission of nuclide $F$} \nonumber \\
    \beta_{k, eff}^{(F)} =& \text{ the effective delayed neutron fraction of \gls{DNP} group $k$ from fission of nuclide $F$. } \nonumber
\end{align}

These equations thus far have not included the energy spectra, but instead have focused on the delayed neutron count rate. 
Equation \eqref{eq:speceq} shows that including the energy in the delayed neutron count rate requires the energy spectra, $\chi_{d, k}(E)$, of the relevant group to be included.
This is the general form of the saturation and pulse irradiation equations, where $A_k$ represents $\bar{\nu}_{d, k} \lambda_k$ for a pulse irradiation or $\bar{\nu}_{d, k}$ for saturation irradiation.
$\tilde{F}_s$ is the total number of fissions for a pulse irradiation or the fission rate for saturation irradiation.
The generic energy and time dependent neutron count rate is given as:

\begin{equation}
\dot{n}_d(E, t) \approx \epsilon \tilde{F}_s \sum_{k=1}^{K} \chi_{d, k} (E) A_k e^{-\lambda_k t},
\label{eq:speceq}
\end{equation}
where, 
\begin{align}
    \tilde{F}_s =& \text{ the generic fission term [- or $s^{-1}$] } \nonumber \\
    A_k =& \text{ the generic matrix term [$s^{-1}$ or -]  } \nonumber \\
    \chi_{d, k}(E) =& \text{ the delayed neutron energy spectrum of \gls{DNP} group $k$ [-]. } \nonumber
\end{align}

\subsubsection{Calculation of DNP Group Spectra and Concentrations}
\label{sec:groupspectra}
Another consideration in calculating group parameters is constructing group concentrations from individual nuclide data.
This is important for situations where reducing the problem from several hundred \glspl{DNP} to a few \gls{DNP} groups reduces computational cost.
A specific use case for this is explicitly modeling all \glspl{DNP} spatially and then using those spatially resolved concentrations with an existing six-group structure.
In such a case, the \glspl{DNP} must be combined in such a manner to preserve their contributions to each \gls{DNP} group.
The reverse is also true, where a given six-group spatially resolved concentration can be separated into individual \glspl{DNP} spatially resolved concentrations.

Two different methods are considered for separating and combining the \gls{DNP} groups into individual nuclides and vice versa.
These methods have been used to generate \gls{DNP} group spectra, so the group spectra can also be calculated using these methods:
 
\begin{itemize}
    \item Grouping based on half lives \cite{england1983aggregate, rudstam1982six}
    \item Grouping based on fractional contribution to adjacent groups \cite{brady1988evaluation}
\end{itemize}

The grouping based on half lives requires comparing each nuclide's half life to a constant value.
The groupings based on half lives are given in the top two sets in Table \ref{tab:halfgroup}.
The bottom set of half lives is from ENDF/B-VIII.0 and is used with the fractional contribution approach \cite{Brown20181}.


\renewcommand{\arraystretch}{1.25}
\begin{table}[H]
    \centering
    \begin{threeparttable}
    \caption{Half-lives of \gls{DNP} groups from various sources.}
    \begin{tabular}{ccccccc}
    \hline
       Source  & $T_1$ $[s]$ &  $T_2$ $[s]$ & $T_3$ $[s]$ & $T_4$ $[s]$ & $T_5$ $[s]$ & $T_6$ $[s]$\\
       \hline
        \cite{england1983aggregate} & $T > 31$ & $31 \geq T > 9$ &  $9 \geq T > 3$ & $3 \geq T > 0.8$ & $0.8 \geq T > 0.3$ & $0.3 \geq T$ \\
        \cite{rudstam1982six}\tnote{*} & $T > 30$ & $30 \geq T > 10$ & $10 \geq T > 4$ & $4 \geq T > 1.4$ & $1.4 \geq T > 0.4$ & $0.4 \geq T$\\
        \cite{Brown20181} & $T = 51.98$ & $T = 21.17$ & $T = 5.74$ & $T = 2.29$ & $T = 0.82$ & $T = 0.24$\\

%        \hline
%        \cite{Brown20181} & $T_1 = 51.98$\\
%        \cite{Brown20181} & $T_2 = 21.17$\\
%        \cite{Brown20181} & $T_3 = 5.74$\\
%        \cite{Brown20181} & $T_4 = 2.29$\\
%        \cite{Brown20181} & $T_5 = 0.82$\\
%        \cite{Brown20181} & $T_6 = 0.24$\\
        \hline
    \end{tabular}
        \begin{tablenotes}
      \item[*] Rudstam uses $T_1 = 55.6$, but to ensure all \glspl{DNP} are included, this has been modified.
    \end{tablenotes}
    \label{tab:halfgroup}
\end{threeparttable}
\end{table}
\renewcommand{\arraystretch}{1}

The half-life grouping method thus results in a series of comparisons; if a nuclide is within a given group's bounds, the concentration of that nuclide is added to that group.

The fractional grouping is more complex, but assigns nuclides to groups in a more rigorous manner.
The method uses two groups, $i$ and $i+1$, for each nuclide $k$.
These groups are determined via Equation \eqref{eq:fraclam} as:
\begin{equation}
    \lambda_i < \lambda_k < \lambda_{i+1}
    \label{eq:fraclam}
\end{equation}
Nuclides which don't fit in these bounds will have the entire fractional contribution go the closest group.
Otherwise, the nuclide will have some fraction of its concentration or spectra go to each \gls{DNP} group.
The fractional contribution, or weight, of nuclide $k$ to each group is represented with $f_{k, i}$ and $f_{k, i+1}$. 
The total fraction must be equal to 1 as shown in Equation \eqref{eq:fracsum} as:
\begin{equation}
    f_{k, i} + f_{k, i+1} = 1
    \label{eq:fracsum}
\end{equation}

To solve for the fractions, the least squares minimization problem is solved as:
\begin{equation}
    \int_{0}^{\infty} \left( \lambda_k e^{-\lambda_k t} -  f_{k, i} \lambda_i e^{-\lambda_i t} - f_{k, i+1} \lambda_{i+1} e^{-\lambda_{i+1} t}   \right)^2 dt
    \label{eq:fracmin}
\end{equation}
Solving this equation for where the partial derivatives are zero for each fraction, $f_{k, i}$ and $f_{k, i+1}$, combined with Equation \eqref{eq:fracsum}, yields the overdetermined system of equations as:
\renewcommand{\arraystretch}{2}
\begin{equation}
A = \begin{bmatrix}
    \lambda_i & \frac{2 \lambda_i \lambda_{i+1}}{\lambda_i + \lambda_{i+1}} \\
    \frac{2 \lambda_i \lambda_{i+1}}{\lambda_i + \lambda_{i+1}} & \lambda_{i+1} \\
    1 & 1\\
\end{bmatrix}, \quad
\vec{b} = \begin{bmatrix}
    \frac{2 \lambda_i \lambda_{k}}{\lambda_i + \lambda_k}\\
    \frac{2 \lambda_k \lambda_{i+1}}{\lambda_k + \lambda_{i+1}} \\
    1\\
\end{bmatrix}, \quad
\vec{x} = \begin{bmatrix}
    f_{k, i} \\
    f_{k, i+1} \\
\end{bmatrix}.
\label{eq:fracmats}
\end{equation}
\renewcommand{\arraystretch}{1}

Because the number of rows is greater than the number of columns in the matrix, the left-inverse is used. This system of equations is then solved as shown in Equation \eqref{eq:lls} as:

\begin{equation}
\vec{x} = (A^T A)^{-1} A^T\vec{b}
\label{eq:lls}
\end{equation}

With the fractional contribution terms, $\vec{x}$, the concentration of a given nuclide can then be multiplied by each fraction to show how much it contributes to that group.
With a summation of each nuclides contribution, the group concentrations can thus be calculated and vice versa.
Alternatively, the spectra can be calculated by using the same approach but with individual \gls{DNP} spectra instead of concentrations.


\subsubsection{Calculation of Kinetics Parameters}
\label{sec:calc-kinetic-param}

The other relevant parameters to be calculated are simpler and do not require solving least squares problems.
The $k^{th}$ \gls{DNP} group delayed neutron fraction is calculated as:
\begin{equation}
    \beta_k = \frac{\bar{\nu}_{d, k}}{\bar{\nu}}.
    \label{eq:betai}
\end{equation}

A useful term for comparing \gls{DNP} group parameters is the average half-life \cite{piksaikin2002energy, piksaikin2007relative}.
Parish et al. also use the average half-life, but do not include the natural log of two in their equation, which could be a typo \cite{parish1999status}.
Equation \eqref{eq:avg-hl} gives the average half-life as:
\begin{equation}
    <T> = \ln (2) \sum_k^{K} \frac{\bar{\nu}_{d, k}}{\lambda_k}.
    \label{eq:avg-hl}
\end{equation}

Equation \eqref{eq:dnyield} shows how the total delayed neutron yield, or average number of delayed neutrons per fission event, $\bar{\nu}_d$, can be calculated \cite{gremyachkin2015verification}:
\begin{equation}
    \bar{\nu}_d = \sum_k^K \bar{\nu}_{d, k}.
    \label{eq:dnyield}    
\end{equation}

This yield can also be approximately represented for an individual nuclide by integrating Equation \eqref{eq:dnpsingleeq} over all time, where the concentration over time is assumed to be: 
\begin{equation}
    N_{i}(t) = N_i(0) e^{-\lambda_i t}.
    \label{eq:conc-time}
\end{equation}
This yields Equation \eqref{eq:delnu-net-nuc}, where the total number of delayed neutrons emitted from a single nuclide $i$ is given as:
\begin{equation}
    \int_0^{\infty} \dot{n}_{d, i}(t) dt= P_{n, i} N_{i} (0).
    \label{eq:delnu-net-nuc}
\end{equation}

The delayed neutron fraction can be written as:
\begin{equation}
    \beta = \frac{\bar{\nu}_d}{\bar{\nu}}
    \label{eq:beta-finalform}
\end{equation}

The mean neutron generation time can be calculated as:
\begin{equation}
    \Lambda = \frac{l}{k},
    \label{eq:lambda-sorta}
\end{equation}
where,
\begin{align}
    l =& \text{ the prompt neutron lifetime [s]} \nonumber \\
    k =& \text{ the neutron multiplication factor [-].} \nonumber \\
\end{align}
This can be calculated using a time-dependent version of a Monte Carlo code along with \gls{PRK} \cite{romero2022calculation}.



Finally, there are the effective parameters: the effective delayed neutron fraction, $\beta_{eff}$, and the effective mean neutron generation time, $\Lambda_{eff}$.
$\beta_{eff}$ is the spectrum and adjoint weighted ratio of the delayed neutron induced fission to the induced fissions from all neutrons \cite{romero2022calculation}.
This is given in Equation \eqref{eq:betaeff-exact} as:

\begin{equation}
    \beta_{eff} = \frac{\int \Phi^\dagger \chi_d \beta \bar{\nu} \Sigma_f \Phi \; dE d\vec{\Omega} d\vec{r}}{\int \Phi^\dagger \chi \bar{\nu} \Sigma_f \Phi \; dE d\vec{\Omega} d\vec{r}},
    \label{eq:betaeff-exact}
\end{equation}
where,
\begin{align}
\Phi^\dagger &= \mbox{ the adjoint neutron flux [$\frac{n}{cm^2 s}$]} \nonumber \\ 
\chi_d &= \mbox{ the delayed neutron energy spectrum [-]} \nonumber \\ 
\chi &= \mbox{ the neutron energy spectrum [-]} \nonumber \\ 
\beta &= \mbox{ the delayed neutron fraction [-]} \nonumber \\ 
\bar{\nu} &= \mbox{ the average number of neutrons released per fission [-]} \nonumber \\ 
\Sigma_f &= \mbox{ the macroscopic fission cross section [$cm^{-1}$]} \nonumber \\
\Phi &= \mbox{ the neutron flux [$\frac{n}{cm^2 s}$].} \nonumber
\end{align}


%\begin{description}
%    %\setlength\itemsep{-0.25em}
%    \item $\phi^\dagger$ is the adjoint neutron flux
%    \item $\chi_d$ is the delayed neutron energy spectrum
%    \item $\chi$ is the fission neutron energy spectrum
%    \item $\beta$ is the delayed neutron fraction
%    \item $\nu$ is the average number of neutrons released per fission
%    \item $\Sigma_f$ is the macroscopic fission cross section
%    \item $\phi$ is the neutron flux
%\end{description}

$\Lambda_{eff}$ can also be calculated explicitly by its definition, which is the "...time between the birth of a neutron and the subsequent absorption that induces fission."
This is shown in Equation \eqref{eq:lambdaeff-exact} as \cite{verboomen2006monte}:
\begin{equation}
    \Lambda_{eff} = \frac{\int \Phi_0^\dagger \frac{1}{\bar{v}} \Phi_0 \; dE d\vec{\Omega} d\vec{r}}{\int \Phi_0^\dagger F_0 \Phi_0 \; dE d\vec{\Omega} d\vec{r}}
    \label{eq:lambdaeff-exact}
\end{equation}
where,
\begin{align}
\Phi_0^\dagger &= \mbox{ the critical adjoint neutron flux [$\frac{n}{cm^2 s}$] } \nonumber \\ 
\Phi_0 &= \mbox{ the critical neutron flux [$\frac{n}{cm^2 s}$] } \nonumber \\ 
\bar{v} &= \mbox{ the average neutron speed [$\frac{cm}{s}$]} \nonumber \\ 
F_0 &= \mbox{ the total fission operator [$cm^{-1}$].} \nonumber
\end{align}

%\begin{description}
%    %\setlength\itemsep{-0.25em}
%    \item $\phi_0^\dagger$ is the critical adjoint neutron flux
%    \item $\phi_0$ is the critical neutron flux
%    \item $v$ is the neutron speed
%    \item $F_0$ is the total fission operator
%\end{description}



There are other methods used to calculate $\beta_{eff}$ and $\Lambda_{eff}$ in the literature \cite{verboomen2006monte}. These include the direct integration, or adjoint
weighted, method \cite{keepin1965physics}; the eigenvalue, or perturbation, method \cite{henry1958application}; and the iterated fission probability method \cite{bohl1964computation}.
$\Lambda_{eff}$ can also be calculated using the pulsed method \cite{romero2022calculation}.
%There is also the pulsed method for calculating $\Lambda_{eff}$, and the prompt method for calculating $\beta_{eff}$ \cite{romero2022calculation}.
For $\beta_{eff}$, there are also the prompt, Meulekamp, and Spriggs methods which can be used \cite{romero2022calculation, meulekamp2006calculating}.
Aufiero et al. have conducted a comparison of the Meulekamp, iterated fission probability, and prompt methods, which show that the iterated fission probability and prompt methods both yield similar results while the Meulekamp method is 200 pcm higher for the \gls{MSFR} \cite{aufiero2014calculating}.